\documentclass[12pt,a4paper]{article}

\usepackage[a4paper,margin=2cm]{geometry}
\usepackage[T2A]{fontenc}
\usepackage[utf8]{inputenc}
\usepackage[russian]{babel}
\usepackage{amsmath,amssymb,mathtools}
\usepackage{array,booktabs,longtable}
\usepackage{xcolor}
\usepackage{hyperref}
\usepackage{tikz}

\hypersetup{
  colorlinks=true,
  linkcolor=blue!55!black,
  urlcolor=blue!55!black
}

\setlength{\parindent}{0pt}
\setlength{\parskip}{0.7em}
\renewcommand{\arraystretch}{1.25}

\begin{document}

\begin{titlepage}
\centering
\vspace*{0.27\textheight}
{\LARGE\bfseries Ревью домашней работы\par}
\vspace{2.2cm}
{\large Автор: Лёвин Артём Александрович\par}
\vspace{0.7cm}
{\large 5 октября 2026 года\par}
\vfill
\begin{tikzpicture}[x=1cm,y=1cm]
  \draw[blue!45!black, line width=0.8pt]
    (0,0) .. controls (1.8,0.55) and (3.6,-0.55) .. (5.4,0)
          .. controls (7.2,0.55) and (9.0,-0.55) .. (10.8,0);
\end{tikzpicture}
\end{titlepage}

\newpage
\section*{Сводная таблица}
\small
\setlength{\tabcolsep}{3pt}
\begin{center}
\begin{tabular}{>{\centering\arraybackslash}p{0.08\linewidth} >{\raggedright\arraybackslash}p{0.20\linewidth} >{\raggedright\arraybackslash}p{0.29\linewidth} >{\raggedright\arraybackslash}p{0.15\linewidth} >{\raggedright\arraybackslash}p{0.15\linewidth}}
\toprule
№ задания & Тема & Суть ошибки & Ваш ответ & Правильный ответ \\
\midrule
\hyperref[task:1]{1} &
Формула двойного угла; область допустимых значений &
После верного шага $2\cos 2x=2$ аргумент функции был изменён: записано $\cos x=1$ вместо $\cos 2x=1$. &
Нет корней. &
Решений нет. \\
\bottomrule
\end{tabular}
\end{center}
\normalsize

\newpage
\vspace*{\fill}
\begin{center}
{\LARGE\bfseries Разбор ошибок}
\end{center}
\vspace*{\fill}
\newpage

\phantomsection\label{task:1}
\section*{Задание 1}

\subsection*{Текст задания}
Решите уравнение
\[
\frac{\sin 4x}{\sin 2x}=2.
\]

\subsection*{Ваше решение}
Сначала было отмечено ограничение
\[
\sin 2x\ne 0.
\]
Затем использована формула двойного угла
\[
\sin 4x=2\sin 2x\cos 2x,
\]
после чего получено
\[
\frac{2\sin 2x\cos 2x}{\sin 2x}=2,
\qquad
2\cos 2x=2.
\]
Далее в записи появляется переход
\[
\cos x=1,
\]
после чего найдено $x=2\pi n$, $n\in\mathbb Z$, и сделан вывод, что подходящих корней нет.

\subsection*{Ваш ответ}
Нет корней.

\subsection*{Ошибка}
\textcolor{red}{Ошибка:} потерян множитель $2$ в аргументе косинуса.

\subsection*{Где именно возникает ошибка}
Последний правильный шаг:
\[
2\cos 2x=2.
\]
Первый неправильный шаг:
\[
\cos x=1.
\]

После деления обеих частей равенства $2\cos 2x=2$ на число $2$ меняется только коэффициент перед косинусом. Аргумент косинуса не меняется. Поэтому должно получиться
\[
\cos 2x=1,
\]
а не $\cos x=1$.

\subsection*{Почему это неверно}
Запись $\cos 2x$ означает, что косинус берётся от аргумента $2x$. Число $2$ перед $x$ находится внутри функции и является частью её аргумента. При делении равенства на $2$ нельзя одновременно изменять аргумент функции.

Иными словами, из равенства
\[
2\cos 2x=2
\]
мы делим левую и правую части на один и тот же внешний множитель $2$:
\[
\frac{2\cos 2x}{2}=\frac{2}{2}.
\]
Получаем именно
\[
\cos 2x=1.
\]
Если заменить $\cos 2x$ на $\cos x$, получится уже другое уравнение с другим множеством кандидатов на корни.

В данном задании конечный вывод «корней нет» всё же совпал с правильным ответом, но это произошло случайно: неверный промежуточный переход дал только часть кандидатов, которые тоже не входят в область допустимых значений. Правильность конечной фразы поэтому не исправляет ошибку в ходе решения.

\subsection*{Ключевая идея}
\textcolor{blue!60!black}{Ключевая идея:} сначала нужно сохранить исходное ограничение знаменателя, затем применить формулу двойного угла и при сокращении помнить, что аргумент $2x$ внутри косинуса остаётся неизменным.

\subsection*{Исправление от последнего правильного шага}
Продолжаем с последнего верного равенства:
\[
2\cos 2x=2.
\]
Делим обе части на $2$:
\[
\cos 2x=1.
\]
Косинус равен единице, когда его аргумент равен целому числу полных оборотов:
\[
2x=2\pi n,\qquad n\in\mathbb Z.
\]
Отсюда
\[
x=\pi n,\qquad n\in\mathbb Z.
\]
Но эти значения ещё обязательно нужно проверить по ограничению исходной дроби.

\subsection*{Полное правильное решение}
В исходном уравнении знаменатель не должен обращаться в нуль:
\[
\sin 2x\ne 0.
\]
Следовательно,
\[
2x\ne \pi k,\qquad k\in\mathbb Z,
\]
то есть
\[
x\ne \frac{\pi k}{2},\qquad k\in\mathbb Z.
\]

Используем формулу
\[
\sin 4x=2\sin 2x\cos 2x.
\]
Тогда исходное уравнение принимает вид
\[
\frac{2\sin 2x\cos 2x}{\sin 2x}=2.
\]
Поскольку по области допустимых значений $\sin 2x\ne 0$, сокращение на $\sin 2x$ допустимо:
\[
2\cos 2x=2.
\]
Делим обе части на $2$:
\[
\cos 2x=1.
\]
Отсюда
\[
2x=2\pi n,\qquad n\in\mathbb Z,
\]
поэтому
\[
x=\pi n,\qquad n\in\mathbb Z.
\]

Теперь проверяем эти кандидаты в ограничении. Для любого целого $n$
\[
\sin 2x=\sin(2\pi n)=0.
\]
Значит, при каждом найденном значении знаменатель исходной дроби равен нулю. Все кандидаты исключаются.

Следовательно, исходное уравнение не имеет решений.

\subsection*{Проверка}
Кандидаты, полученные после преобразований, имеют вид $x=\pi n$. При каждом таком значении знаменатель исходного выражения равен нулю:
\[
\sin(2\pi n)=0.
\]
Поэтому ни один кандидат нельзя подставить в исходную дробь. Допустимых корней нет.

\subsection*{Итог}
\textcolor{teal!60!black}{Итог:} правильный ответ --- решений нет. Ошибка находится в переходе от $2\cos 2x=2$ к уравнению с косинусом от $x$: аргумент $2x$ нельзя менять при делении на внешний коэффициент.

\subsection*{Задача на закрепление}
Решите уравнение
\[
\frac{\sin 6x}{\sin 3x}=\sqrt{2}.
\]

\vfill
\begin{center}
\small Лёвин Артём Александрович эксклюзивно для Ярослава Гаврилова
\end{center}

% Сообщение Ярославу по итогам ревью:
% В первой задаче конечный ответ получился правильным, но в ходе решения потерялся множитель два внутри аргумента косинуса: после равенства «два косинуса от двух икс равно двум» нужно было получить «косинус от двух икс равен единице».

\end{document}
